\subsection*{从执行顺序推导乘法顺序}
先把图形横向拉长两倍，再剪切，怎样用一个矩阵保存整个过程？必须知道每根标准基箭头经过两步后落在哪里。第一步的矩阵 $B$ 已把位置存进自己的列；第二步对每一列作用 $A$。因此复合矩阵的列是 $A$ 乘以 $B$ 的各列，这就是矩阵乘法的来源。

$ABv$ 的括号是 $A(Bv)$：书写从左到右，执行从最靠近输入的右端开始，与 $f(g(x))$ 一样。把乘法读成动作，还能判断不同形状的矩阵能否相乘：前一步的输出维度必须接上后一步的输入维度。

\subsection*{交换顺序为什么改变结果}
拿 $(0,1)$ 做实验。先横向拉长时，它在纵轴上，不动；再剪切，横坐标变成 $1$。若先剪切，它先到 $(1,1)$，随后拉长会把刚产生的横坐标放大，结果是 $(2,1)$。第二步处理的是已经改变的对象，所以同样两步不必得到同一结果。

再看公式里的求和：第 $i$ 个输出分量，要累加中间层各分量对它的贡献。下标 $k$ 标记所有中间通道。神经网络层间计算与换基都会重用这个视角。

三个动作按 $C,B,A$ 执行，先把哪两步打包，不改变实际过程，所以 $(AB)C=A(BC)$。结合律允许重新打包，却没有允许交换动作。以后读长矩阵公式时，给每个矩阵配一句实际指令，再决定能怎样移动括号。
